\section{Methods}
\label{sec:method}

\begin{figure*}[t]
\centering
\includegraphics[width=\textwidth]{fig1_overview.png}
\caption{\textbf{Overview of scRareRefine.} A frozen scANVI backbone produces a latent embedding and a prediction
for every cell. On the training set we build class prototypes; for each query cell we compute an anisotropic
rare-membership score. A validation set calibrates a conformal threshold and three abstention gates so that the
false-rescue rate stays below $\alpha$. At test time, candidate cells whose score exceeds the threshold are
relabelled to the rare class. Reference prototypes use only training cells; all thresholds use only validation cells.}
\label{fig:overview}
\end{figure*}

\subsection{Setup and inductive protocol}
\label{sec:method-setup}

Let a dataset contain cells with expression profiles and a single designated rare type $\rare$. We split cells into
training, validation and test partitions using a batch-/donor-held-out scheme (target $70/15/15$), so that
validation and test cells come from donors unseen during training; this is the harder, more realistic inductive
setting. A semi-supervised scANVI model \citep{xu2021scanvi}, itself initialised from an scVI model
\citep{lopez2018scvi}, is trained on the training partition with all majority cells labelled and only a limited
budget of rare cells labelled. The labelled rare budget is $\max(5, \lfloor p\,\Ntrain^{\rare}\rfloor)$ for a
nominal fraction $p$, where $\Ntrain^{\rare}$ is the number of rare cells in the training partition; the floor of
five reflects the minimum needed to estimate a prototype. The model yields, for every cell, a latent embedding
$\latent\in\Real^{d}$ and a predicted label. scRareRefine consumes these frozen outputs and never retrains the
backbone.

\subsection{Prototype-membership scoring}
\label{sec:method-proto}

For each class $c$ observed among the \emph{labelled} training cells we form a prototype
$\proto{c}=\frac{1}{|S_c|}\sum_{i\in S_c}\latent_i$, the mean latent of its labelled cells $S_c$, and a robust radius
$\radius{c}=\max\!\big(\mathrm{median}_{i\in S_c}\,\lVert \latent_i-\proto{c}\rVert,\,\epsilon\big)$
(set to a conservative default when $|S_c|<3$). For a query cell with embedding $\latent$ we define an
\emph{anisotropic rare-membership score}
\begin{equation}
\score(\latent) \;=\; \softmax_{c}\!\Big(-\frac{\lVert \latent-\proto{c}\rVert}{\radius{c}}\Big)\Big[\rare\Big],
\label{eq:score}
\end{equation}
i.e.\ the rare-class component of a softmax over per-class distances normalised by each class's compactness. The
per-class normalisation lets the score adapt to classes of different spread and limits intrusion from a tight
neighbouring majority cluster. We also define a train-only \emph{separability ratio}
\begin{equation}
\sep \;=\; \frac{\min_{c\neq \rare}\lVert \proto{\rare}-\proto{c}\rVert}{\frac{1}{|S_\rare|}\sum_{i\in S_\rare}\lVert \latent_i-\proto{\rare}\rVert},
\label{eq:sep}
\end{equation}
the distance from the rare prototype to its nearest majority prototype divided by the mean rare intra-class radius;
$\sep$ measures how geometrically distinct the rare type is, using only training information.

\subsection{Conformal calibration and abstention gates}
\label{sec:method-conformal}

Given the frozen scores, scRareRefine decides which test cells to relabel using only validation labels. Three gates
and one threshold are applied in sequence (\cref{alg:rescue}).

\paragraph{Separability gate.} If $\sep<\sep_{\mathrm{low}}$ the rare and majority populations are too entangled for
prototype rescue to be safe and the module abstains. We fix $\sep_{\mathrm{low}}=1.3$ as a dataset-independent prior
(\cref{sec:results-sep}).

\paragraph{Necessity guard.} If the validation backbone already recovers the rare type --- specifically if the number
of validation rare cells it misses is below $M$ --- rescue is unnecessary and may only add noise, so the module
abstains. We fix $M=3$, the minimum support needed for the validation-adaptive selection below to carry signal.

\paragraph{Conformal threshold.} Let the validation non-rare scores be $\score_1,\dots,\score_n$. We set the
threshold $\tautest$ to the finite-sample $(1-\alphabudget)$ conservative order statistic
$\tautest=\score_{(k^\star)}$ with $k^\star=\lceil(1-\alphabudget)(n+1)\rceil$ (and $\tautest=+\infty$, i.e.\ no
rescue, if $k^\star>n$). Under exchangeability of validation and test non-rare cells, relabelling only cells with
$\score\ge\tautest$ bounds the probability that a non-rare cell is rescued at $\alphabudget$.

\paragraph{Validation-adaptive candidate rank.} A cell is a \emph{candidate} only if the rare prototype is among its
$k$ nearest prototypes (isotropic Euclidean rank $\le k$) and its backbone label is not already rare. We choose
$k\in\{1,2,3\}$ on validation as the value that maximises validation rare-cell F1 subject to a Wilson $95\%$ upper
bound on the validation FFR not exceeding $\alphabudget$, breaking ties toward the smaller (more conservative) $k$.
Using the Wilson upper bound rather than a point estimate guards against the validation/test distribution shift
induced by donor-held-out splits.

\paragraph{Rescue rule.} A test cell is relabelled to $\rare$ iff it is a candidate (rank $\le k$) \emph{and}
$\score\ge\tautest$. All other predictions are left unchanged.

\begin{algorithm}[t]
\caption{scRareRefine conformal rescue (inductive)}
\label{alg:rescue}
\begin{algorithmic}[1]
\Require frozen latents and predictions for train/val/test; rare type $\rare$; budget $\alphabudget$
\State fit prototypes $\{\proto{c}\}$, radii $\{\radius{c}\}$, separability $\sep$ on labelled \emph{train} cells
\If{$\sep<\sep_{\mathrm{low}}$} \Return backbone predictions \Comment{separability gate}
\EndIf
\State $v_{\mathrm{miss}}\gets$ \#validation rare cells missed by backbone
\If{$v_{\mathrm{miss}}<M$} \Return backbone predictions \Comment{necessity guard}
\EndIf
\State compute scores \cref{eq:score} for val and test; set $\tautest$ from val non-rare scores
\State choose $k\in\{1,2,3\}$ maximising val rare-F1 s.t.\ Wilson-upper val-FFR $\le\alphabudget$
\State \textbf{for} each test cell: relabel to $\rare$ if rank$\le k$ and $\score\ge\tautest$
\State \Return refined test predictions
\end{algorithmic}
\end{algorithm}

\subsection{False-rescue rate}
\label{sec:method-ffr}

We quantify the cost of rescue by the \emph{false-rescue rate}
$\ffr=\big|\{\text{non-rare test cells relabelled to }\rare\}\big|\,/\,\big|\{\text{non-rare test cells}\}\big|$.
The conformal threshold and gates are designed so that $\ffr$ is controlled at the fixed budget $\alphabudget=0.01$.
Because validation and test cells come from different donors under our split, this control is empirical rather than a
strict finite-sample guarantee; we report $\ffr$ with raw counts and treat $\alphabudget$ as a fixed prior, never
tuned per dataset.
